% =====================================================================
%  PAPER 2  —  A Ruppeiner Curvature Dichotomy for Generalized Prime Gases
%  IG-PRIMON-T1.  FULL assembled draft v0.1 (2026-06-14).
%  Compilable end-to-end: pdflatex twice (for refs/labels).
%  Math core locked [V]; prose drafted. Bose table = computed values
%  (the v0.1 draft's tau=0.1, 0.01 rows were interpolated and are corrected).
%  Citations verified 2026-06-14 (bakas1991, julia1994 titles corrected).
%  For J. Phys. A submission, swap \documentclass to {iopart}.
% =====================================================================
\documentclass[a4paper,11pt]{article}
\usepackage{amsmath,amssymb,amsthm,mathtools}
\usepackage{booktabs}
\usepackage[margin=1in]{geometry}
\usepackage[colorlinks=true,linkcolor=blue,citecolor=blue]{hyperref}

\theoremstyle{plain}
\newtheorem{theorem}{Theorem}
\newtheorem{lemma}{Lemma}
\newtheorem{proposition}{Proposition}
\theoremstyle{definition}
\newtheorem{definition}{Definition}
\theoremstyle{remark}
\newtheorem{remark}{Remark}

\newcommand{\Ric}{R}
\newcommand{\Li}{\operatorname{Li}}
\newcommand{\Hess}{\operatorname{Hess}}
\newcommand{\eps}{\varepsilon}
\newcommand{\detg}{\det g}

\title{A Ruppeiner curvature dichotomy for generalized prime gases:\\
       asymptotic flatness at limiting-temperature transitions}
\author{Aaron Jones\\ \small JtechAI LLC}
\date{\today}

\begin{document}
\maketitle

\begin{abstract}
We study the thermodynamic (Ruppeiner) geometry of grand-canonical \emph{generalized prime gases},
whose partition functions are Dirichlet series $\log\Xi(\beta,z)=\sum_{k\ge1}(z^k/k)\,\mathcal P(k\beta)$.
We prove that the scalar curvature at a phase transition is governed by \emph{where the singularity
lives}: a \emph{temperature-driven} (Hagedorn / limiting-temperature) transition forces the curvature
to vanish asymptotically, $\Ric\to0$ (the theorem), while a \emph{fugacity-driven} (condensation) transition
breaks the underlying cancellation, and we exhibit the resulting divergence $|\Ric|\to\infty$ at
Bose--Einstein condensation. The discriminating mechanism is a single algebraic condition on the
singular sector of $\log\Xi$ --- it is curvature-flat precisely when that sector factorizes as
$e^{y}\varphi(x)$ in the natural coordinates, i.e.\ when fugacity and temperature do not couple. The
result is established by an explicit cancellation lemma together with an analyticity argument, and is
bracketed by four exactly-computable instances: the classical ideal gas ($\Ric=0$), the Riemann
primon gas and an all-integer Bose gas (both Hagedorn, $\Ric\to0$), and the ideal Bose gas at
condensation ($|\Ric|\to\infty$). Two of the four contain no primes; the statement is one of
statistical mechanics.
\end{abstract}

% =====================================================================
\section{Introduction}\label{sec:intro}
The Ruppeiner scalar curvature $R$ of the thermodynamic state space offers a geometric measure of
underlying statistical fluctuations, classically diverging at critical points where the correlation
volume becomes macroscopic \cite{ruppeiner1995}. For one-dimensional manifolds such curvature claims
are vacuous; in a two-parameter (grand-canonical) ensemble, however, $R$ is a non-trivial invariant
\cite{ruppeiner1995,janyszek1990}.

In this paper we investigate the thermodynamic geometry of grand-canonical generalized prime gases
(Beurling gases) near their continuous phase transitions. For the canonical primon gas, the Hagedorn
temperature ($\beta=1$) induces a well-known divergence in the partition function
\cite{julia1990,julia1994,spector1990,bakas1991}. We extend this to a two-parameter grand-canonical
manifold evaluated along the locking ray $z=1$.

We establish a rigorous curvature dichotomy: for generalized prime gases characterized by a definite
leading singularity, the scalar curvature $R$ strictly vanishes (asymptotically flat) at
temperature-driven (Hagedorn-like) singularities, but diverges ($|R|\to\infty$) at fugacity-driven
(Bose--Einstein condensation) transitions. We prove that the flat sector is unconditional within this
class, driven by an exact decoupling in the Hessian determinant.

% =====================================================================
\section{Setup}\label{sec:setup}

\begin{definition}[generalized prime gas]\label{def:gas}
A grand-canonical gas with
\begin{equation}\label{eq:logXi}
   \log\Xi(\beta,z)=\sum_{k\ge1}\frac{z^k}{k}\,\mathcal P(k\beta),
\end{equation}
where $\mathcal P$ is the spectral (prime) zeta of the mode spectrum, $\beta$ the inverse temperature
and $z$ the fugacity.
\end{definition}

Use natural coordinates $(x,y)=(-\beta,\log z)$. The family is exponential with sufficient statistics
$(E,N)$, so the Ruppeiner/Fisher metric is the Hessian of the potential $\psi:=\log\Xi$,
\begin{equation}\label{eq:metric}
   g=\Hess_{(x,y)}\psi,\qquad
   \partial_x^{a}\partial_y^{b}\psi=(-1)^{a}\sum_{k\ge1}k^{\,a+b-1}z^{k}\,\mathcal P^{(a)}(k\beta).
\end{equation}
For a two-dimensional Hessian metric the scalar curvature is
\begin{equation}\label{eq:R}
   \Ric=-\frac{N}{2(\detg)^{2}},\qquad
   N=\det\!\begin{pmatrix}
       \psi_{xx} & \psi_{xy} & \psi_{yy}\\
       \psi_{xxx}& \psi_{xxy}& \psi_{xyy}\\
       \psi_{xxy}& \psi_{xyy}& \psi_{yyy}
     \end{pmatrix}.
\end{equation}

\begin{remark}[convention]\label{rem:convention}
We use the sphere-positive scalar convention, in which the normal family
$\psi=-t_1^2/(4t_2)-\tfrac12\log(-2t_2)$ returns $\Ric=-1$ at every point; \eqref{eq:R} reproduces
this and the flat-product control $\Ric=0$. Only $|\Ric|$ (vanishing vs.\ divergence) is convention-
independent; reported signs are fixed by this choice.
\end{remark}

\paragraph{Hypotheses.}
\begin{itemize}
\item[\textbf{(H1)}] \emph{Temperature-driven singularity.} $\mathcal P(s)$ is real-analytic on
$(1,\infty)$ with an isolated singularity at its abscissa $s=1$ and no other singularity on the real
axis $s\ge1$. (The transition is a singularity in $\beta$ at the limiting temperature $\beta=1$.)
\item[\textbf{(H2)}] \emph{Finite background.} $\sum_{k\ge2}k(k-1)\mathcal P^{(j)}(k)<\infty$ for
$j=0,1,2$ (automatic when $\mathcal P(k)\to0$ geometrically).
\item[\textbf{(C1)}] \emph{Definite leading order.} The abscissa singularity is of pole type
$A\eps^{-\alpha}$ ($\alpha>0$) or logarithmic. (Holds for every natural Beurling/prime spectral zeta.)
\item[\textbf{(C2)}] \emph{Nondegenerate amplitude.} $\Delta_3:=\sum_{k\ge2}k(k-1)\mathcal P(k)\neq0$
(automatic when the spectral coefficients are positive, $\mathcal P(k)>0$). Used only to fix the explicit
amplitude in Theorem~\ref{thm:flat}; flatness itself does not require it (Remark~\ref{rem:uncond}).
\end{itemize}

% =====================================================================
\section{The product-form criterion}\label{sec:lemmas}

\begin{lemma}[product form annihilates the curvature numerator]\label{lem:C1}
If $\psi_{\mathrm{sing}}=e^{y}\varphi(x)$ for a smooth $\varphi$, then the third row of $N$ in
\eqref{eq:R}, $(\psi_{xxy},\psi_{xyy},\psi_{yyy})$, equals the first row
$(\psi_{xx},\psi_{xy},\psi_{yy})$ identically and at all orders; hence the contribution of
$\psi_{\mathrm{sing}}$ to $N$ is $\equiv0$.
\end{lemma}
\begin{proof}
$\partial_y\!\big(e^{y}\varphi^{(j)}(x)\big)=e^{y}\varphi^{(j)}(x)$, so each entry of row three equals
the corresponding entry of row one. A determinant with two equal rows vanishes; the identity is exact
in $\eps$, hence order by order.
\end{proof}

\begin{lemma}[only the single-particle term is singular]\label{lem:k1}
Under \textnormal{(H1)}, every $k\ge2$ term of \eqref{eq:logXi} is real-analytic at
$(\beta,z)=(1,1)$, and the unique singular term is
$\psi_1=z\,\mathcal P(\beta)=e^{y}\varphi(x)$, $\varphi(x):=\mathcal P(-x)$.
\end{lemma}
\begin{proof}
For $k\ge2$ and $\beta$ near $1$, $\mathcal P(k\beta)$ is evaluated at $k\beta\approx k\ge2$, inside
the analytic region by (H1); the term is analytic. Only $k=1$ probes the abscissa, and that term is
exactly $z\,\mathcal P(\beta)$---linear in $z$, hence literally of product form $e^{y}\varphi(x)$ with
$x=-\beta$, $y=\log z$, $\varphi(x)=\mathcal P(-x)$.
\end{proof}

% =====================================================================
\section{The dichotomy theorem}\label{sec:thm}

\begin{theorem}[temperature-driven $\Rightarrow$ asymptotically flat]\label{thm:flat}
Under \textnormal{(H1)--(H2)} and \textnormal{(C1)}, as $(\beta,z)\to(1,1)$ along $z=1$,
$\Ric\to0$. Explicitly, with $L=\log(1/\eps)$, $\eps=\beta-1$, background amplitude
$\Delta_3:=\sum_{k\ge2}k(k-1)z^{k}\mathcal P(k)$, and leading singularity $\mathcal P(1+\eps)\sim A\eps^{-\alpha}$,
\[
   \alpha>0:\ \Ric\sim\frac{\Delta_3(\alpha+1)}{2A^{2}}\,\eps^{2\alpha},
   \qquad
   \alpha\to0\ (\text{logarithmic}):\ \Ric\sim\frac{\Delta_3}{2(L+\kappa)^{2}}.
\]
\end{theorem}
\begin{proof}
By Lemma~\ref{lem:k1} the singular sector is $\psi_1=e^{y}\varphi(x)$; by Lemma~\ref{lem:C1} it
contributes $0$ to $N$, so $N$ is sourced entirely by the analytic $k\ge2$ background. The punchline is a
rate mismatch: the surviving (background-sourced) numerator scales as $N\sim\eps^{-2\alpha-4}$ while the
metric volume diverges faster, $(\detg)^2\sim\eps^{-4\alpha-4}$, so $\Ric=-N/2(\detg)^2\sim\eps^{2\alpha}\to0$.
The full cofactor reduction and the resulting rates are given in Appendix~\ref{app:rates}; the flatness
conclusion is unconditional within the class (Remark~\ref{rem:uncond}).
\end{proof}

\begin{proposition}[fugacity-driven $\Rightarrow$ cancellation fails]\label{prop:div}
If the transition is at $z\to z_c$ with the singularity carried by the full fugacity series,
$\log\Xi_{\mathrm{sing}}=\varphi(\beta)\,\Li_a(z)$, then the singular sector
$\varphi(\beta)\,\Li_a(e^{y})$ is not of the form $e^{y}\varphi(x)$: Lemma~\ref{lem:C1} does not apply,
$\mathrm{Row}_3\ne\mathrm{Row}_1$, and the product-form cancellation \emph{fails}, so $\Ric$ is no longer
forced to vanish. The resulting divergence $|\Ric|\to\infty$ is realized at Bose--Einstein condensation
(\S\ref{sec:bose}, Table~\ref{tab:bose_divergence}).
\end{proposition}

\begin{remark}\label{rem:criterion}
Theorem~\ref{thm:flat} and Proposition~\ref{prop:div} give the discriminant in one line:
\[
   \boxed{\ \Ric\to0 \iff \text{the singular sector of }\log\Xi\text{ factorizes as } e^{y}\varphi(x).\ }
\]
\end{remark}

% =====================================================================
\section{Instances}\label{sec:instances}
All four instances run through the same engine \eqref{eq:R}; the classical ideal gas validates it
($\Ric=0$ is forced and obtained).

\begin{table}[h]
\centering
\begin{tabular}{@{}llll@{}}
\toprule
gas & singular structure & type & curvature (rate / amplitude)\\
\midrule
classical ideal gas        & $u^{-3/2}e^{y}\ (\varphi\cdot e^{y})$ & product form     & $\Ric\equiv0$ (exact, product form)\\
Riemann primon gas         & $z\,P(\beta),\ P\sim\log(1/\eps)$    & temp-driven, log & $\Ric\to0^{+}$, $\sim1/(L+\kappa)^2$, $\Delta_3=5.0451882$\\
all-integer Bose gas       & $z(\zeta(\beta)-1)\sim1/\eps$        & temp-driven, pole& $\Ric\to0^{+}$, $\Ric/\eps^2\to\Delta_3=5.693982$\\
ideal Bose gas (BEC)       & $u^{-3/2}\Li_{5/2}(e^{y})$           & fugacity-driven  & $|\Ric|\to\infty$, $\sim\tau^{-1/2}$\\
\bottomrule
\end{tabular}
\caption{The dichotomy, bracketed. $u=\beta$; $\tau=-\log z\to0^{+}$ at condensation;
$\kappa=0.832503$ for the primon gas. The two flat instances differ in spectrum (primes vs.\ all
integers) and singularity type (log vs.\ pole); two of the four contain no primes.}
\label{tab:instances}
\end{table}

\paragraph{All-integer Bose gas.} $\prod_{n\ge2}(1-z\,n^{-\beta})^{-1}$, i.e.\ $\mathcal P=\zeta-1$,
along $z=1$: $\Ric/\eps^2 = 4.95,\,5.61,\,5.69$ at $\eps=10^{-2},10^{-3},10^{-4}$, converging to
$\Delta_3^{\mathrm{int}}=5.693982$.

\subsection{The fugacity-driven complement: ideal Bose gas}\label{sec:bose}
To exhibit the divergent branch we evaluate the ideal Bose gas, where the singularity is
fugacity-driven: as $\tau\to0$ the macroscopic occupation of the ground state controls the
fluctuations. Table~\ref{tab:bose_divergence} reports the computed scalar curvature on approach to
BEC. The divergence $|\Ric|\to\infty$ is unambiguous; the product $\Ric\,\tau^{1/2}$ drifts slowly
($0.18\to0.11$ as $\tau:0.3\to0.001$) toward a finite limit, consistent with the expected BEC leading
behaviour $|\Ric|\sim\tau^{-1/2}$ up to the standard subleading corrections (the finite-$\tau$ effective
exponent approaches $\tfrac12$ from below).

\begin{table}[h]
\centering
\renewcommand{\arraystretch}{1.15}
\begin{tabular}{c|cc}
$\tau=-\log z$ & $\Ric$ & $\Ric\,\tau^{1/2}$\\
\hline
$0.3$   & $0.3275$ & $0.1794$\\
$0.1$   & $0.4681$ & $0.1480$\\
$0.03$  & $0.7511$ & $0.1301$\\
$0.01$  & $1.2135$ & $0.1214$\\
$0.003$ & $2.1196$ & $0.1161$\\
$0.001$ & $3.5874$ & $0.1134$\\
\end{tabular}
\caption{Scalar curvature of the ideal Bose gas approaching BEC ($z=e^{-\tau}$, sphere-positive
convention): $|\Ric|\to\infty$, leading $\sim\tau^{-1/2}$.}
\label{tab:bose_divergence}
\end{table}

\begin{remark}[a geometric invariant of the primon gas]\label{rem:Cconst}
The renormalized information length of the Riemann primon gas from the Hagedorn point to $\beta=2$,
\[
   C=\int_{1}^{2}\!\Big[\sqrt{\partial_\beta^2\log\zeta(\beta)}-\tfrac{1}{\beta-1}\Big]\,d\beta
    =-0.0343561541791219860831108814584476\ldots,
\]
is, to $90$ digits, not an integer combination of $\{\pi,\log2,\log3,\gamma,\gamma_1,\zeta(3),\zeta(5),
G,\log A\}$ at coefficient height $\le10^{6}$ (PSLQ); we register it as a new constant of the gas.
\end{remark}

% =====================================================================
\section{Discussion}\label{sec:disc}
The dichotomy demonstrates that the thermodynamic geometry of arithmetic gases discriminates strictly
by the source of the singularity. Because the singular sector of a temperature-driven transition
factorizes as $e^{y}\varphi(x)$, it generates an exactly degenerate (flat) contribution to the
Hessian determinant; the surviving curvature is sourced entirely by the finite analytic background
coupling to that singular mode, and $\Ric\to0$.

Conversely, fugacity-driven transitions avoid this cofactor cancellation, permitting the standard
Ruppeiner divergence $|\Ric|\to\infty$. The general rate amplitude
$\Ric\sim[\Delta_3(\alpha+1)/(2A^2)]\eps^{2\alpha}$ (Appendix~\ref{app:rates}) shows the flatness
holds across arbitrary pole orders $\alpha>0$ and the marginal logarithmic limit. We stress the
precise status: the flatness is \emph{unconditional} within the class --- it survives even the
degenerate background $\Delta_3=0$, with $\Ric$ vanishing one power faster ($\eps^{2\alpha+1}$); the
positive spectral condition $\Delta_3>0$, automatic whenever $\mathcal P(k)=\sum_q q^{-k}>0$, fixes
only the explicit rate \emph{amplitude}. Thus the Hagedorn transitions of arithmetic gases are not
merely physically distinct from ordinary condensation, but geometrically flat --- and this is a
statement of statistical mechanics: two of the four bracketing instances (the classical ideal gas and
the all-integer Bose gas) involve no primes, and nothing here bears on the location of any zeta zero.

% =====================================================================
\appendix
\section{Curvature rates: the cofactor reduction}\label{app:rates}

Evaluate on $z=1$, $\beta=1+\eps$, $\eps\to0^{+}$, and abbreviate the singular (single-particle)
derivatives $p_j:=\mathcal P^{(j)}(1+\eps)$, $j=0,1,2,3$.

\subsection*{A.1 Singular/background split and the row reduction}
By Lemma~\ref{lem:k1} the $k=1$ term is the only singular one; write each partial as its $k=1$
(singular) part plus the finite $k\ge2$ background. From \eqref{eq:metric} at $z=1$ the singular parts
are $\psi_{xx}^{s}=p_2$, $\psi_{xy}^{s}=-p_1$, $\psi_{yy}^{s}=p_0$, $\psi_{xxx}^{s}=-p_3$,
$\psi_{xxy}^{s}=p_2$, $\psi_{xyy}^{s}=-p_1$, $\psi_{yyy}^{s}=p_0$, so the three rows of $N$ have
singular parts
\[
   \mathrm{Row}_1^{s}=(p_2,-p_1,p_0),\quad \mathrm{Row}_2^{s}=(-p_3,p_2,-p_1),\quad
   \mathrm{Row}_3^{s}=(p_2,-p_1,p_0)=\mathrm{Row}_1^{s}.
\]
The equality $\mathrm{Row}_3^{s}=\mathrm{Row}_1^{s}$ is Lemma~\ref{lem:C1} in coordinates: the $k=1$
weight $k^{a+b-1}=1$ is independent of $b$. Subtracting Row~1 from Row~3 annihilates the divergent
singular part, leaving the finite background differences
\begin{equation}\label{eq:rowdiff}
   \mathrm{Row}_3-\mathrm{Row}_1=(\Delta_{xx},\Delta_{xy},\Delta_3),\quad
   \Delta_{xx}=\!\sum_{k\ge2}\!k(k-1)\mathcal P''(k),\
   \Delta_{xy}=-\!\sum_{k\ge2}\!k(k-1)\mathcal P'(k),\
   \Delta_3=\!\sum_{k\ge2}\!k(k-1)\mathcal P(k).
\end{equation}
Determinant-invariance under $\mathrm{Row}_3\to\mathrm{Row}_3-\mathrm{Row}_1$ gives
\begin{equation}\label{eq:Ncofactor}
   N=\Delta_{xx}M_{31}-\Delta_{xy}M_{32}+\Delta_3M_{33},
\end{equation}
with minors on rows $1,2$ whose leading (singular$\times$singular) parts are
\begin{equation}\label{eq:minorlead}
   M_{31}\sim p_1^2-p_0p_2,\qquad M_{32}\sim p_0p_3-p_1p_2,\qquad M_{33}\sim p_2^2-p_1p_3,
\end{equation}
and $\detg\sim p_0p_2-p_1^2$.

\subsection*{A.2 Pole singularity of order $\alpha$}
For $\mathcal P(1+\eps)\sim A\eps^{-\alpha}$, $\alpha>0$, $A\ne0$:
$p_0\sim A\eps^{-\alpha}$, $p_1\sim-A\alpha\eps^{-\alpha-1}$, $p_2\sim A\alpha(\alpha+1)\eps^{-\alpha-2}$,
$p_3\sim-A\alpha(\alpha+1)(\alpha+2)\eps^{-\alpha-3}$. Substituting,
\[
   \detg\sim A^2\alpha\,\eps^{-2\alpha-2},\qquad M_{33}\sim-A^2\alpha^2(\alpha+1)\,\eps^{-2\alpha-4},
\]
and since the minor orders are $M_{31}\sim\eps^{-2\alpha-2}$, $M_{32}\sim\eps^{-2\alpha-3}$,
$M_{33}\sim\eps^{-2\alpha-4}$, the $M_{33}$ term dominates \eqref{eq:Ncofactor}:
$N\sim\Delta_3M_{33}=-\Delta_3A^2\alpha^2(\alpha+1)\eps^{-2\alpha-4}$. Hence
\begin{equation}\label{eq:poleRate}
   \Ric=-\frac{N}{2(\detg)^2}
   \sim\frac{\Delta_3A^2\alpha^2(\alpha+1)\eps^{-2\alpha-4}}{2A^4\alpha^2\eps^{-4\alpha-4}}
   =\frac{\Delta_3(\alpha+1)}{2A^2}\,\eps^{2\alpha}\ \xrightarrow[\eps\to0]{}\ 0 .
\end{equation}
The singular powers cancel between $N$ and $(\detg)^2$; the surviving $\eps^{2\alpha}$ forces flatness.

\subsection*{A.3 Logarithmic singularity (marginal $\alpha\to0$)}
For $\mathcal P(1+\eps)=\log(1/\eps)+a_0+\cdots$, $L=\log(1/\eps)$: $p_0\sim L$, $p_1\sim-\eps^{-1}$,
$p_2\sim\eps^{-2}$, $p_3\sim-2\eps^{-3}$, and the cancellations of A.2 are replaced by
logarithmically-enhanced leading terms, $\detg\sim(L+\kappa)/\eps^2$, $M_{33}\sim-\eps^{-4}$,
$N\sim-\Delta_3/\eps^4$, where $\kappa$ collects the $O(1)$ singular-subleading and
singular$\times$background contributions ($\kappa=0.832503$ for the primon gas). Hence
\begin{equation}\label{eq:logRate}
   \Ric\sim\frac{\Delta_3}{2(L+\kappa)^2}\ \xrightarrow[\eps\to0]{}\ 0,
\end{equation}
the $\alpha\to0^{+}$ boundary of \eqref{eq:poleRate}: $\eps^{2\alpha}\to1$ ceases to vanish and the
logarithmic growth of $\detg$ supplies the vanishing instead.

\subsection*{A.4 Conditions and the rate claim}
Two hypotheses were used. (C1) (definite leading order) holds for any Beurling/prime spectral zeta.
(C2) $\Delta_3\ne0$ is needed only for the explicit amplitude; since
$\Delta_3=\sum_{k\ge2}k(k-1)\mathcal P(k)$ with $\mathcal P(k)=\sum_q q^{-k}>0$, $\Delta_3>0$
automatically for every physical gas.

\begin{remark}[flatness is unconditional]\label{rem:uncond}
If $\Delta_3=0$ (non-physical), the $M_{33}$ term drops and $N$ is led by
$\Delta_{xy}M_{32}\sim\eps^{-2\alpha-3}$, giving $\Ric\sim\eps^{2\alpha+1}\to0$ --- one power faster.
Thus $\Ric\to0$ for \emph{every} temperature-driven singularity in the class, independent of (C2);
only the amplitude depends on $\Delta_3$.
\end{remark}

\subsection*{A.5 Numerical confirmation}
Eq.~\eqref{eq:poleRate} predicts $\Ric/\eps^{2\alpha}\to\Delta_3(\alpha+1)/(2A^2)$. On the tunable-order
family $\mathcal P(s)=(s-1)^{-\alpha}2^{1-s}$ ($A=1$):
\[
\renewcommand{\arraystretch}{1.15}
\begin{array}{c|ccc|c}
\alpha & \eps=10^{-2} & 10^{-3} & 10^{-4} & \text{predicted}\\\hline
1 & 2.708 & 2.968 & 2.997 & 3.000\\
2 & 2.402 & 2.526 & 2.538 & 2.540\\
3 & 2.366 & 2.531 & 2.549 & 2.551
\end{array}
\]
The marginal case \eqref{eq:logRate} is confirmed by the primon gas ($\Ric(L+\kappa)^2\to\Delta_3/2=2.5226$)
and the $\alpha=1$ case by the all-integer Bose gas ($\Ric/\eps^2\to\Delta_3=5.6940$).

% =====================================================================
\begin{thebibliography}{9}
\bibitem{bakas1991}
I. Bakas and M. J. Bowick, ``Curiosities of arithmetic gases,'' \textit{J. Math. Phys.} \textbf{32}(7),
1881--1884 (1991).
\bibitem{janyszek1990}
H. Janyszek and R. Mruga\l{}a, ``Riemannian geometry and stability of ideal quantum gases,''
\textit{J. Phys. A: Math. Gen.} \textbf{23}, 467 (1990).
\bibitem{julia1990}
B. Julia, ``Statistical theory of numbers,'' in \textit{Number Theory and Physics}, eds.\ J.-M. Luck,
P. Moussa, M. Waldschmidt, Springer Proc.\ Phys.\ \textbf{47}, 276--293 (1990).
\bibitem{julia1994}
B. L. Julia, ``Thermodynamic limit in number theory: Riemann-Beurling gases,'' \textit{Physica A}
\textbf{203}, 425 (1994).
\bibitem{ruppeiner1995}
G. Ruppeiner, ``Riemannian geometry in thermodynamic fluctuation theory,'' \textit{Rev. Mod. Phys.}
\textbf{67}, 605 (1995).
\bibitem{spector1990}
D. Spector, ``Supersymmetry and the M\"obius inversion function,'' \textit{Commun. Math. Phys.}
\textbf{127}, 239 (1990).
\end{thebibliography}

\end{document}
